\paragraph{AQ-136-PATH-LIFT: the arithmetic path arrow need not admit source lifting.}
Let $P=\{p_1,\ldots,p_s\}$ be nonempty and finite, and let $k$
be a $\mathbb Z_{(P)}$-algebra with
$K=k\otimes_{\mathbb Z}\mathbb Q\ne0$.
Use the exact arithmetic replacement $\mathcal E$ and natural
arrow $\eta_C:C\to\mathcal E(C)$ of AQ-136-SEMILOCAL.
The following proposed source-lifting property is false:
\[
 \begin{gathered}
 F:A\longrightarrow\mathcal E(C),\quad A,C\in\mathcal A_k
 \quad\Longrightarrow\quad
 \text{there exist }A'\in\mathcal A_k,\ h:A'\to A,\ u:A'\to C\\
 \text{ with }Fh=\eta_Cu.
 \end{gathered}                                             \tag{136.PL1}
\]
In particular it fails for $k=\mathbb Q$ and for mixed semilocal
rings such as $\mathbb Z_{(2,3)}$.
This is a failure of a specified way of refining an arrow;
it does not refute common sources or terminal localization.

\begin{proof}
Use the standard $B$ and its square-zero extension $D=B\ltimes M$
from AQ-136-STRICT. To fix every operation needed here,
\[
 B_n=\Gamma_k(x_0,\ldots,x_n)\otimes_k
       \Lambda_k(e_0,\ldots,e_n)/(\textstyle\sum e_j),
 \quad dx_j=e_j,\quad T=\textstyle\sum x_j .                \tag{136.PL2}
\]
The module $M$ is free over $B$ on $\alpha_j$ of form degree
zero and $\beta_j$ of form degree one, subject to
$\sum\alpha_j=\sum\beta_j=0$, with
$d\alpha_j=\beta_j$ and $d\beta_j=0$.
All four kinds of generator have weight one.
The degree-zero relative special derivation
\[
 \delta:B\longrightarrow M,\qquad
 \delta x_j=\alpha_j,\quad\delta e_j=\beta_j               \tag{136.PL3}
\]
commutes with $d$ and is zero on $S$.
The extension has $M^2=0$ and
$\gamma_r(b+m)=\gamma_r(b)+\gamma_{r-1}(b)m$ on its
positive-weight even ideal. Consequently
\[
       j_0(b)=(b,0),\qquad j_1(b)=(b,\delta b)              \tag{136.PL4}
\]
are strict simplicial PD differential $S$-algebra maps.

The earlier full proof verifies $B,D\in\mathcal A_k$.
Its bounds can also be read directly here: $M$ is vertically
contractible by
$h(b\beta_j)=(-1)^{|b|}b\alpha_j$, $h(b\alpha_j)=0$,
so $H^*(D_n)=S$ in every weight. Horizontally, with
$Q=L/k\sum x_j$ as before,
\[
 M_*^{w,q}=(B_*^{w-1,q}\otimes Q)
                 \oplus(B_*^{w-1,q-1}\otimes Q).
\]
Both summands are contractible for $w>q+1$ by the checked
linear contraction of the standard rows. Thus all bounds
are uniform in the simplex and no coefficient tail is removed.

Define a relative derivation $\iota:B\to M$ of form degree
$-1$ and weight zero by
\[
                  \iota(x_j)=0,\qquad\iota(e_j)=\alpha_j.
                                                               \tag{136.PL5}
\]
Equivalently, $\iota=h\delta$ for the graded $B$-linear
contraction of $M$ displayed above. Extend it by the signed
Leibniz and special divided-power rules. These rules preserve the exterior relations: for
example $\iota(e_j^2)=\alpha_je_j-e_j\alpha_j=0$.
They preserve $\sum e_j=0$ because $\sum\alpha_j=0$.
They also preserve every coefficient, and commute with all
simplicial operators, which are the same fibre sums on the
four types of generator. On generators and divided-power
monomials one has
\[
                           d\iota+\iota d=\delta.          \tag{136.PL6}
\]
Both sides are degree-zero special derivations, so this
identity holds on the whole algebra.

For signs, write the rational de Rham path algebra with the
interval factor first,
$K[z,dz]\otimes_K D_K$; the graded switch identifies it
with the path algebra used in AQ-136-SEMILOCAL. Define
\[
 H:B_K\longrightarrow K[z,dz]\otimes_K D_K,\qquad
 H(b)=1\otimes(b,0)+z\otimes(0,\delta b)
                         +dz\otimes(0,\iota b).           \tag{136.PL7}
\]
This formula is homogeneous in form degree and weight.
The square-zero property of $M$ makes it multiplicative:
the $z$ term is the ordinary derivation rule, while the
$dz$ term is the signed rule for $\iota$, together with
$(-1)^{|a|}$ when $1\otimes a$ crosses $dz$.
The unit and coefficients are fixed. Divided powers are
preserved by the same first-order square-zero formula
and the special derivation identities; equivalently, over
$K$ they are uniquely the operations $v\mapsto v^r/r!$.

It is a chain map. For homogeneous $b$, its module component
has differential
\[
 z\otimes\delta(db)+dz\otimes(\delta b-d\iota b)
          =z\otimes\delta(db)+dz\otimes\iota(db),
\]
by (136.PL6), which is the module component of $H(db)$.
The formula is simplicially natural by the generator check
above. Finally its endpoints are $j_{0,K}$ and $j_{1,K}$.

We now obtain an arrow over the original coefficient ring:
\[
 F:B\longrightarrow\mathcal E(D),\qquad
 b\longmapsto
 \bigl((j_0(b_i))_{i=1}^s,\ j_1(b_K),\ (H(b_K))_{i=1}^s\bigr),
                                                               \tag{136.PL8}
\]
where $b_i$ and $b_K$ are the scalar extensions of $b$.
The endpoints in (136.PL7) are exactly the compatibility
conditions defining $\mathcal E(D)$. Thus (136.PL8) is a
strict morphism in the original category.

Suppose (136.PL1) held for this arrow. Every path component
of $Fh=\eta_Du$ would be constant. Evaluating either endpoint
therefore gives
\[
                  j_{0,K}h_K=j_{1,K}h_K,
                    \qquad\text{hence}\quad\delta h_K=0.  \tag{136.PL9}
\]
The kernel of $\delta:B_K\to M_K$ is exactly $S_K$ in every
simplex. Indeed, eliminate $x_0,e_0$ using the coefficient
and sum relations. In the resulting basis
$x_1^{[a_1]}\cdots x_n^{[a_n]}e_I$ over $S_K$, the
$\beta_j$ coefficient of $\delta$ first eliminates every
monomial containing $e_j$. Then its $\alpha_j$ coefficient,
which lowers $a_j>0$ by one with coefficient one, eliminates
every positive $a_j$. The remaining terms are precisely
$S_K$. This argument uses freeness of the displayed basis,
not any field hypothesis on $K$.

Consequently $h_K:A'_K\to B_K$ factors through the constant
simplicial algebra $S_K$. The factorization is an $S_K$-map,
so it retracts $S_K\to A'_K$. In each weight it makes the
nonzero constant module $KT^{[w]}$ a direct summand of the
zero-form row of $A'_K$. That row cannot be acyclic in the
weight tail. On the other hand $A'_K$ is admissible: flat
localization preserves both uniform vertical cohomology
and all normalized horizontal homology. This contradiction
proves failure of (136.PL1).
\end{proof}

Over a nonzero $\mathbb Q$-algebra, common sources and ordinary
terminal localization were separately proved in AQ-136-RATIONAL.
Thus this counterexample cannot be promoted to a negative answer
to either original question. It rules out the general source-lifting
shortcut for the actual natural arithmetic path arrow, including
its mixed semilocal instances; a more restricted lifting assertion
would require a separate proof and narrower hypotheses.
